
\subsubsection{Key Findings}

The phase transition represents restructuring, not fragmentation. CV and NLP dynamics were never unified (r = $-0.13$ pre-2020). Foundation models were associated with attention reallocation from traditional to emergent benchmarks without fragmenting coordinated behavior---because none existed.

The mechanism chain is: foundation models emerge $\rightarrow$ new benchmarks created ($19\times$ acceleration) $\rightarrow$ researcher attention shifts (+19\% to emergent benchmarks) $\rightarrow$ traditional benchmarks persist with reduced dominance.

\subsubsection{Unexpected Finding}

The refutation of h-m5 yielded a novel insight: modality-specific benchmark usage patterns were independent before 2020, challenging the implicit assumption of unified pre-transition dynamics in phase transition models. The slight positive correlation post-2021 (r = +0.23) could indicate that foundation models, which often bridge CV and NLP, may have introduced more cross-modality coherence.

\subsubsection{Practical Implications}

The attention shift has concrete consequences for practitioners. ImageNet, despite commanding 47,068 papers, now captures only 11.70\% of community attention, while emergent benchmarks capture over 26\%. This gap shapes visibility, fundability, and perceived research relevance. Teams achieving SOTA on traditional benchmarks may find their work framed as incremental, while comparable performance improvements on emergent benchmarks signal engagement with foundation model capabilities.

\subsubsection{Limitations}

\begin{itemize}
\item \textbf{Causal inference:} The analysis demonstrates temporal coincidence between foundation model emergence and benchmark concentration changes. While the mechanism verification chain provides supporting evidence, strict causation cannot be established. Confounding factors (field growth, funding shifts, publication venue dynamics) may contribute.
\end{itemize}

\begin{itemize}
\item \textbf{h-m1 data source:} Citation counts were sourced from Google Scholar due to API timeout. Z-scores are extreme enough that variations would not change the verdict, but exact values are approximate.
\end{itemize}

\begin{itemize}
\item \textbf{PWC coverage:} Papers With Code may not capture industry or proprietary research.
\end{itemize}

\begin{itemize}
\item \textbf{Temporal resolution:} Monthly aggregation is approximate to $\pm 1$--2 months.
\end{itemize}

\begin{itemize}
\item \textbf{Single baseline:} Only a monotonic trend null hypothesis was tested. Robustness to alternative change-point methods (Bayesian, CUSUM) and PELT parameter sensitivity remains untested.
\end{itemize}

\begin{itemize}
\item \textbf{Modality classification:} Sparse coverage in some modalities (Audio, Tabular) limits reliability of cross-modality correlation analysis.
\end{itemize}
